%! Populations, Samples, Parameters, and Statistics — Unit 1, Lesson 1.2 — Exit Ticket (Answer Key)
\documentclass[10pt]{article}
\usepackage{saar-article}
\usepackage{saar-key}   % pulls in -boxes; do NOT also load -boxes
\newcommand{\TallMath}[1]{$\displaystyle #1\rule[-1.4em]{0pt}{3.2em}$}

\begin{document}

\pageheader{Unit 1, Lesson 1.2}{Exit Ticket --- Answer Key}
% NAMESTRIP: no \namedateperiod here — mirrors the blank. Teacher-only prose goes
% in the lesson plan as \begin{teachernote}[Exit Ticket], never in a key.

\vspace{0.15in}

\begin{notesbox}{Exit Ticket --- No Notes}
\small
Westfield High School has $800$ students. The school newspaper asked $40$ students
in one lunch period where they eat lunch. Of those $40$, $26$ said the cafeteria
--- that is $65\%$.

\begin{enumerate}[leftmargin=1.6em, itemsep=10pt, topsep=3pt]

  \item Name the groups in this study.

        Population: \ans{all 800 Westfield students} \hfill Sample size: \ans{40}

        Sample: \ans{the 40 students asked in one lunch period}

  \item Label each number \textbf{parameter} or \textbf{statistic}.

        \begin{center}
        \renewcommand{\arraystretch}{1.45}
        \begin{tabularx}{0.95\linewidth}{@{} X c @{}}
        \toprule
        \textbf{The number} & \textbf{Which kind?} \\
        \midrule
        $800$ students attend Westfield.                        & \ans{parameter} \\
        $65\%$ of the $40$ students asked eat in the cafeteria. & \ans{statistic} \\
        The true percent of all $800$ Westfield students who eat in the cafeteria.
          & \ans{parameter} \\
        \bottomrule
        \end{tabularx}
        \end{center}

  \item \textbf{(a)} Name one constraint that kept the newspaper from asking all
        $800$ students. \ans{time --- one lunch period is not enough}

        \textbf{(b)} Tomorrow the newspaper asks a \emph{different} $40$ students
        and gets $60\%$. Does that mean one of the two surveys is wrong? Explain.
        \ansline{No. A statistic changes from sample to sample --- that is sampling variability.}
        \ansline{Both estimate the same one true percent for all 800 students.}

\end{enumerate}
\end{notesbox}

\end{document}
